\documentclass[11pt,a4paper]{article}
\usepackage[italian]{babel}
\usepackage[margin=2.5cm]{geometry}
\usepackage{parskip}
\usepackage{amsmath, amssymb, amsthm}
\usepackage{booktabs}
\usepackage{array}
\usepackage{xcolor}
\usepackage{needspace}
\usepackage{enumitem}
\usepackage{listings}
\usepackage{tikz}
\usepackage[most]{tcolorbox}
\usepackage[hidelinks]{hyperref}
\usetikzlibrary{decorations.pathreplacing, shapes.geometric, arrows.meta, positioning}

% Niente vedove/orfane: i paragrafi non lasciano righe isolate a fine/inizio pagina
\widowpenalty=10000
\clubpenalty=10000

% Elenchi più compatti
\setlist{itemsep=2pt, parsep=0pt, topsep=4pt, partopsep=0pt}

% --- Riquadri con bordo nero, senza numero (si spezzano tra due pagine) ---
% Uso: \begin{definizione} ... \end{definizione}
%      \begin{teorema}[Nome] ... \end{teorema}  (il nome è facoltativo)
\tcbset{nota/.style={enhanced jigsaw, breakable, colback=white,
  colframe=black, boxrule=0.5pt, sharp corners}}
\NewTColorBox{definizione}{}{nota,
  before upper={\textbf{Definizione.}\ }}
\NewTColorBox{teorema}{o}{nota,
  before upper={\textbf{Teorema\IfValueT{#1}{ (#1)}.}\ }}
\NewTColorBox{proposizione}{o}{nota,
  before upper={\textbf{Proposizione\IfValueT{#1}{ (#1)}.}\ }}
\NewTColorBox{proprieta}{o}{nota,
  before upper={\textbf{Proprietà\IfValueT{#1}{ (#1)}.}\ }}
\NewTColorBox{assioma}{o}{nota, boxrule=1pt,
  before upper={\textbf{Assioma\IfValueT{#1}{ (#1)}.}\ }}

% --- Ambienti semplici (senza riquadro) ---
% Il titolo sta in cima al blocco, anche se il contenuto inizia con un riquadro o
% con codice; e non resta mai da solo in fondo a una pagina.
% Uso: \begin{esempio}[Titolo facoltativo] ... \end{esempio}
\newcommand{\newrunin}[2]{%
  \NewDocumentEnvironment{#1}{o}{%
    \par\Needspace{4\baselineskip}\addvspace{\medskipamount}\noindent
    \textbf{#2}\IfValueT{##1}{ (##1)}\textbf{.}\ \ignorespaces}%
  {\par\addvspace{\medskipamount}}}
\newrunin{esempio}{Esempio}
\newrunin{esercizio}{Esercizio}
\newrunin{osservazione}{Osservazione}

% V e F nelle tavole di verità
\newcommand{\V}{\textsf{V}}
\newcommand{\F}{\textsf{F}}

% --- Codice C ---
% Uso: \begin{codice} ... \end{codice}      codice C numerato
%      \begin{comando} ... \end{comando}    comandi da terminale
%      \begin{risultato} ... \end{risultato} output di un programma
%      \lstinline|printf("ciao");|           codice dentro al testo
\lstdefinestyle{c}{
  language=C,
  basicstyle=\ttfamily\small,
  keywordstyle=\color{blue}\bfseries,
  commentstyle=\color{gray}\itshape,
  stringstyle=\color{red!70!black},
  numbers=left,
  numberstyle=\tiny\color{gray},
  numbersep=8pt,
  columns=flexible,
  keepspaces=true,
  breaklines=true,
  showstringspaces=false,
  tabsize=4,
  morekeywords={include, define},
  % lettere accentate nei commenti
  literate={à}{{\`a}}1 {è}{{\`e}}1 {é}{{\'e}}1 {ì}{{\`i}}1
           {ò}{{\`o}}1 {ù}{{\`u}}1 {È}{{\`E}}1 {À}{{\`A}}1
}
\lstset{style=c}
\newtcblisting{codice}{listing only, nota, left=6mm,
  listing options={style=c}}
\newtcblisting{comando}{listing only, nota,
  listing options={basicstyle=\ttfamily\small, numbers=none, language={},
    columns=flexible, keepspaces=true, breaklines=true}}
\newtcblisting{risultato}{listing only, enhanced jigsaw, breakable,
  colback=black!5, colframe=black, boxrule=0.5pt, sharp corners,
  listing options={basicstyle=\ttfamily\small, numbers=none, language={},
    columns=flexible, keepspaces=true, breaklines=true}}

% --- Diagrammi di flusso (simboli standard) ---
\tikzset{
  terminale/.style = {ellipse, draw, minimum width=2.5cm, minimum height=0.9cm},
  io/.style        = {trapezium, trapezium left angle=70, trapezium right angle=110,
                      draw, minimum width=2.5cm, minimum height=0.9cm},
  processo/.style  = {rectangle, draw, minimum width=2.5cm, minimum height=0.9cm},
  decisione/.style = {diamond, draw, aspect=2, minimum width=2.5cm},
  freccia/.style   = {thick, -{Stealth}}
}

\title{Calcolo combinatorio: regole di base e notazioni}
\author{Marco Cerbella}
\date{Lezione 7 -- 12 ottobre 2026}

\begin{document}
\maketitle

\section{Regole di somma e di prodotto}

Siano $A, B$ insiemi finiti con $|A| = n$ e $|B| = k$.

\begin{proprieta}[Regola della somma]
Se $A \cap B = \emptyset$ allora $|A \cup B| = n + k$.
\end{proprieta}

\begin{proprieta}[Regola del prodotto]
$|A \times B| = n \cdot k$.
\end{proprieta}

Infatti gli elementi di $A \times B$ sono le coppie $(a_i, b_j)$ e si possono disporre in una tabella con $n$ righe (una per ogni $a_i$) e $k$ colonne (una per ogni $b_j$):
\begin{center}
\begin{tabular}{c|ccc}
$A \times B$ & $b_1$ & $\cdots$ & $b_k$ \\ \midrule
$a_1$ & $(a_1, b_1)$ & $\cdots$ & $(a_1, b_k)$ \\
$\vdots$ & $\vdots$ & & $\vdots$ \\
$a_n$ & $(a_n, b_1)$ & $\cdots$ & $(a_n, b_k)$
\end{tabular}
\end{center}

\begin{esempio}
Una collana ha una pietra, che può essere Blu, Verde o Rossa, e un simbolo, che può essere Cuore, Stella, Cerchio, Quadrato o Fiore. Si possono costruire $3 \cdot 5 = 15$ collane con una pietra e un simbolo.
\end{esempio}

\section{Notazioni}

\begin{definizione}
Il \emph{fattoriale} di $n \in \mathbb{N}$ è
\[
n! = \begin{cases} 1 & \text{se } n = 0, \\ n \cdot (n-1)! & \text{se } n \geq 1, \end{cases}
\]
cioè $n! = n \cdot (n-1) \cdots 2 \cdot 1$.
\end{definizione}

\begin{definizione}
Per $n \geq k$ il \emph{coefficiente binomiale} ``$n$ su $k$'' è
\[
\binom{n}{k} = \frac{n!}{k!\,(n-k)!} = C_{n,k}.
\]
\end{definizione}

\section{Come si scelgono $k$ oggetti tra $n$}

Dati $n$ oggetti, ne scegliamo $k$: in quanti modi possiamo fare la scelta? La risposta dipende da due domande: l'\emph{ordine} conta? sono ammesse \emph{ripetizioni}? Prendiamo come esempio $X = \{x, y, z\}$.

\begin{center}
\renewcommand{\arraystretch}{1.4}
\begin{tabular}{@{}l|p{5.2cm}|p{5.2cm}@{}}
 & \textbf{Insiemi} (l'ordine non conta) & \textbf{Liste} (l'ordine è importante) \\ \midrule
\textbf{Senza ripetizioni}
 & \textbf{Combinazioni semplici}: $\{x,y\} = \{y,x\}$; $\{x,x\}$ non accettabile
 & \textbf{Disposizioni semplici}: $(x,y) \neq (y,x)$; $(x,x)$ non accettabile \\ \midrule
\textbf{Con ripetizioni}
 & \textbf{Combinazioni con ripetizioni}: $\{x,y\} = \{y,x\}$; $\{x,x\}$ accettabile
 & \textbf{Disposizioni con ripetizioni}: $(x,y) \neq (y,x)$; $(x,x)$ accettabile
\end{tabular}
\end{center}

\end{document}
